\section{ChokeBench: Grading Localisers on the Target That Matters}\label{sec:chokebench}
ProcessBench grades a verifier by whether it finds the human-labelled first error. The value curves
make a second target measurable: the \emph{choke point}, the step where $v$ falls most. ChokeBench
grades step localisers on both, and on what the named step is worth as a resampling point. It uses
the \cbN\ fatal traces whose outcome is still open ($v_0>\delta$; effective $n=\cbNeff$ after sample
splitting) and the released step signals; no model is re-run.

\textbf{Protocol.} A localiser names one step $j$ of a failed trace. We report (i) whether $j$ is
the human first error; (ii) the posterior probability that $j$ is the choke point (Jeffreys
posterior on every $v_k$, 2,000 draws), so sampling noise is not counted as a miss; (iii) the
per-attempt success of resampling from just before $j$, minus restart; and (iv) the same difference
in P(solved) at a FLOP-matched budget of four restart attempts, \emph{charging nothing for finding
$j$}. Item (iv) is an upper bound on any verifier that uses the localiser; Table~\ref{tab:policies}
charges the verifier. Half of each prefix's continuations select the open traces (and the oracle's
choke point), the other half score them, over 50 splits; the estimator reproduces the per-model
values of Table~\ref{tab:backtrack} within sampling error.

\begin{table}[t]\centering\small
\caption{ChokeBench, pooled over four generators. Accuracy on the human first error (Human) and on
the choke point (Choke, posterior); success minus restart per attempt and at a 4$\times$ budget
with free localisation. 95\% trace-bootstrap CIs in Fig.~\ref{fig:chokebench}.}
\label{tab:chokebench}
\setlength{\tabcolsep}{3pt}
\begin{tabular}{lcccc}\toprule
Localiser & Human & Choke & $\Delta$/attempt & $\Delta$ at 4$\times$\\\midrule
Human first error & 1.00 & \cbHumanAccD & \cbHumanGain & \cbHumanGainB\\
PRM: first step $<0.5$ & \cbPrmFirstAccH & \cbPrmFirstAccD & \cbPrmFirstGain & \cbPrmFirstGainB\\
PRM: largest drop & \cbPrmDropAccH & \cbPrmDropAccD & \cbPrmDropGain & \cbPrmDropGainB\\
PRM: minimum & \cbPrmMinAccH & \cbPrmMinAccD & \cbPrmMinGain & \cbPrmMinGainB\\
Generator: max token NLL & \cbNllAccH & \cbNllAccD & \cbNllGain & \cbNllGainB\\
Generator: max entropy & \cbEntAccH & \cbEntAccD & \cbEntGain & \cbEntGainB\\
Random step & \cbRandAccH & \cbRandAccD & \cbRandGain & \cbRandGainB\\
First step (= restart) & \cbFirstAccH & \cbFirstAccD & \cbFirstGain & \cbFirstGainB\\\midrule
Learned, tuned (held-out model) & \cbLearnedAccH & \cbLearnedAccD & \cbLearnedGain & \cbLearnedGainB\\
Oracle choke point (split) & -- & -- & \cbOracleGain & \cbOracleGainB\\\bottomrule
\end{tabular}
\end{table}

\begin{figure*}[t]\centering
\includegraphics[width=0.9\textwidth]{fig25_chokebench.pdf}
\caption{ChokeBench. Left: PRM rules find human-labelled errors far more often than choke points.
Right: P(solved) minus restart at a budget of four restart attempts when finding the step is free.}
\label{fig:chokebench}
\end{figure*}

\textbf{Localisers find errors, not choke points} (Table~\ref{tab:chokebench},
Fig.~\ref{fig:chokebench}). The PRM's largest-drop rule names the human first error in
\cbPrmDropAccHPct\% of open failures, consistent with its ProcessBench F1, but the choke point in
only \cbPrmDropAccDPct\% [\cbPrmDropAccDLo, \cbPrmDropAccDHi]. The expert label itself is the choke
point with probability \cbHumanAccD. Generator uncertainty (maximum token NLL or entropy) is far
below the PRM on both targets, though above a random step. The beginning of the trace is the choke point
in \cbFirstAccDPct\% of open failures: many of them go wrong on the first step.

\textbf{Even a perfect localiser buys little.} Per attempt, every deployable localiser is worse than
restart, from \cbPrmFirstGain\ to \cbEntGain, and so is the human label (\cbHumanGain). With the
step found for free and the shorter continuations converted into extra attempts, the human label
is still \cbHumanGainB\ [\cbHumanGainBLo, \cbHumanGainBHi] and the best PRM rule
\cbPrmFirstGainB. An oracle that knows the choke point from held-out rollouts reaches
\cbOracleGainB\ [\cbOracleGainBLo, \cbOracleGainBHi]. That is the ceiling for any backtracking
verifier on these traces, before its own cost is paid. It is small because, before the choke point,
a failed trace is worth about as much as a fresh start (Appendix~\ref{app:prop}): finding the choke
point saves tokens, not successes.

\textbf{Can a localiser be learned?} We train one directly for the choke-point target. Each step
of a failed trace becomes an example described by nine features (PRM score, PRM drop, first PRM
score below 0.5, mean and maximum token NLL and entropy, step length, relative position), labelled
softly with the posterior probability that it is the choke point. The model names the step with the
highest predicted probability. Evaluation is leave-one-generator-out, with nested selection: for
each held-out generator, \ltGrid\ configurations (logistic regression with $C\in\{0.01,0.1,1,10\}$ and
gradient-boosted trees with depth $\in\{2,3\}$ and learning rate $\in\{0.05,0.2\}$, each on PRM,
generator or all features) are scored by leave-one-generator-out cross-validation on the other three
generators only, and the best is refit and tested once (Fig.~\ref{fig:tuning}). The tuned localiser
finds the choke point with probability \cbLearnedAccD\ [\cbLearnedAccDLo, \cbLearnedAccDHi], no better
than the fixed PRM rules, and its restart point is worth \cbLearnedGainB\ at a 4$\times$ budget.
Inner-CV scores do not carry over to the held-out model (best inner score \ltQwenMathInner\ vs test
\ltQwenMathTest\ for Qwen2.5-Math-7B): which step decides a failure is model-specific, as the
transfer failure of the learned restart policy in Section~\ref{sec:policies} already suggested.

\begin{figure*}[t]\centering
\includegraphics[width=0.85\textwidth]{fig26_learned_tuning.pdf}
\caption{Nested hyperparameter selection of the learned localiser. Each line is one outer fold
(held-out generator); points are inner leave-one-generator-out scores of the 24 configurations;
stars mark the selected configuration; the legend gives its score on the held-out generator.}
\label{fig:tuning}
\end{figure*}

\textbf{Release.} \texttt{results/chokebench/chokebench\_v1.jsonl} holds, for every failed trace,
the per-prefix success counts, continuation counts and lengths, the human label and the PRM scores.
A new localiser is graded by naming one step per trace (\texttt{scripts/49\_chokebench.py}).
